\section{Return lifts, physical observations, and phase coverage}
\label{sec:v83-geometry}

The scalar inverse recovers absolute functions, not merely their differentials.
This distinction matters for return time. We formulate the geometric
consequence on the domains on which the recovered functions are actual
generating sheets, and then exhibit the freedom left by an unobserved
open phase region.

Let $X\subset\R^{2h}$ have canonical one-form $\lambda=p\,dx$.
An exact lift is a pair $g=(F,A)$ satisfying $dA=F^*\lambda-\lambda$.
Composition and inversion are
\begin{equation}\label{eq:v83-lifts}
 (F,A)(G,B)=(F\circ G,B+A\circ G),\qquad
 (F,A)^{-1}=(F^{-1},-A\circ F^{-1}).
\end{equation}
An absolute type-I generating sheet $S(x,y)$ represents the graph
\begin{equation}\label{eq:v83-generating}
 p=-S_x(x,y),\qquad p'=S_y(x,y),\qquad A(x,p)=S(x,y).
\end{equation}
It is regular when $S_{xy}$ is invertible. No minimizing property or
convexity of $S$ is required.

\begin{theorem}[Selective records determine a covered return lift]
\label{thm:v83-return}
Let $r,s$ be coprime positive integers. For each of $g^r,g^s$, suppose
finitely many regular endpoint gates in fixed canonical coordinates
are given. On each gate, the actual absolute generating sheets are
the components of Theorem~\ref{thm:v83-main}, with a fixed detector
degree $q$. The sheet count, channels, actions, and pairing between
gates are not supplied. Suppose their phase graphs cover all domains
required by a fixed word $(g^r)^u(g^s)^v$, where $ur+vs=1$,
including the domains of inverse factors. Assume fixed collars of
these domains and that the sheets belong to the indicated powers of
one deterministic lift.

Then $q+3$ projective deadlines on each gate determine $g$ on the
covered target. Under the margins of Theorem~\ref{thm:v83-global},
fixed $C^{k+2}$ action bounds, and uniform inverse mixed-Jacobian and
phase-map inverse bounds, there is a locally Lipschitz reconstruction
from $C^{k+2}$ normalized matrix data to the $C^k$ lift.
\end{theorem}
\begin{proof}
Theorem~\ref{thm:v83-main} gives the unordered absolute sheets on each
gate, including any action coincidences. Their differentials in
\eqref{eq:v83-generating} determine the phase graphs. At an initial
phase point, regularity solves $p=-S_x(x,y)$ locally for $y$.
Two solutions in overlapping charts have the same output and the
same primitive: by hypothesis they represent one actual deterministic
lift. Thus equality of initial phase coordinates pairs the recovered
graphs, without supplying semantic names to their component indices.
This recovers $g^r,g^s$ on the domains in the word. Since they are
powers of the same lift, \eqref{eq:v83-lifts} gives
\[
                         g=(g^r)^u(g^s)^v.
\]
Every factor is evaluable by the coverage assumption. In particular
the absolute primitive constants are retained.

We specify the local noisy construction as well. Fix a finite cover
of the reference phase target by smaller charts whose closures lie
inside the implicit-function domains, together with a subordinate
partition of unity. On each chart select the reference sheet and the
local solution of $p=-S_x$; Theorem~\ref{thm:v83-global} identifies
its perturbed continuation. Uniform mixed-Jacobian bounds give
$C^k$ errors for the local map and primitive from $C^{k+2}$ sheet
errors. The extra action derivative supplies bounded $C^{k+1}$
phase maps. On overlaps these exact maps agree in the \emph{fixed
canonical coordinates}. Therefore one may average their nearby
coordinate-valued maps and primitives with the fixed partition. This
step is not an average of unmatched sheet indices. For inverse factors
use, on the fixed inverse-domain cover, the local inverse branches
continued from the reference map; their estimates follow from the
inverse-Jacobian bound. Compose the resulting factors in the specified
word on its nested collars. Product, chain, and inverse estimates give
the asserted $C^k$ bound for this finite construction. The extension
need not be exactly symplectic for data outside the model, nor do its
independently extended forward and inverse factors have to be exact
inverses off the model. Both properties hold on exact data, which is
all that the uniqueness assertion requires.
\end{proof}

\begin{corollary}[An existing transverse return section]
\label{cor:v83-contact}
For a smooth Reeb flow with a smooth transverse global section, let
$F$ and $\tau>0$ be the first-return map and time. If the actual
return sheets of two coprime counts satisfy Theorem~\ref{thm:v83-return},
selective endpoint records recover $(F,\tau)$ and its section-preserving
strict contact suspension on the covered target.
\end{corollary}
\begin{proof}
Write $\lambda$ for the restriction of the contact form to the section.
The section flow-box map pulls back the contact form to $dt+\lambda$.
Differentiating the variable-time return identity gives
$F^*\lambda-\lambda=d\tau$. Thus $(F,\tau)$ is an exact lift and
its powers have the actual elapsed return times as primitives. Apply
Theorem~\ref{thm:v83-return}; the quotient
$(x,t)\sim(Fx,t-\tau(x))$ preserves $dt+\lambda$.
\end{proof}

The following separation of supplied and recovered information is part
of the statement, not an auxiliary observation protocol.
\begin{center}
\small
\begin{tabular}{@{}p{.42\textwidth}p{.49\textwidth}@{}}
\toprule
Supplied geometric or apparatus structure & Recovered from the matrices \\
\midrule
Endpoint gates and canonical coordinates & Visible sheet count and local component covering \\
Return counts and their common deterministic lift & Absolute action functions on the gates \\
Phase cover, inverse-factor domains, and collars & Canonical graphs and their coordinate pairing \\
Regularity and fixed finite detector degree & Relative detector coefficients and weights \\
Two conditionally independent clock-invariant, full-rank readouts & Both unknown readout channels \\
An existing transverse global section, when a flow is considered & Return map, return time, and the covered suspension \\
\bottomrule
\end{tabular}
\end{center}
The theorem does not infer an unobserved global section or continue
through an unobserved caustic. The next result explains why some global
acquisition hypothesis is unavoidable in the unrestricted smooth class.

\subsection{An exact ambiguity outside the sampled orbit tube}

\begin{theorem}[An unsampled open phase region]
\label{thm:v83-unobserved}
Let $g=(F,A)$ be a smooth exact lift on a canonical phase domain.
For finitely many positive integers $n$, suppose even the full lifted
restrictions $g^n|_{D_n}$ are known, where the $D_n$ have compact
closures and all relevant iterates are defined. Set
\[
       K=\bigcup_n\bigcup_{j=1}^n F^j(\overline{D_n}).
\]
If an open set $O\Subset X\setminus K$ satisfies
$F^{-1}(O)\cap D_*\ne\varnothing$ for an open target $D_*$,
there are arbitrarily small smooth exact perturbations of $g$ which
have identical lifted observations on every $D_n$ but differ on $D_*$.
The perturbation may be chosen to preserve any strict positive-roof
margin on a fixed compact neighborhood of the target.
\end{theorem}
\begin{proof}
Choose $z_*\in F^{-1}(O)\cap D_*$. A smooth Hamiltonian $H$
compactly supported in $O$, with nonzero Hamiltonian vector field at
$F(z_*)$, has a small-time flow $h_\epsilon$ which is the identity
outside $O$ and moves $F(z_*)$ for all sufficiently small nonzero
$\epsilon$. With the convention
$\iota_{X_H}d\lambda=-dH$, Cartan's formula gives the compactly
supported primitive
\[
 B_\epsilon=\int_0^\epsilon
           (\lambda(X_H)-H)\circ h_t\dd t,
 \qquad h_\epsilon^*\lambda-\lambda=dB_\epsilon.
\]
Both $h_\epsilon-\id$ and $B_\epsilon$ vanish outside $O$.
Define $g_\epsilon=(h_\epsilon,B_\epsilon)g$.
For $z\in D_n$, induction on $j\le n$ shows that
$g_\epsilon^j(z)=g^j(z)$ as lifted maps: each successive original
phase image $F^j(z)$ lies in $K$, where the added lift is the identity
with zero primitive. The induction also shows that the accumulated
primitive is unchanged, not merely the phase map. At $z_*$ the maps
differ. Smooth compact support gives $g_\epsilon\to g$ in every
fixed finite smooth norm on the relevant compact domains. A strict
roof margin therefore persists.
\end{proof}

This is a necessity result for smooth dynamics with restricted orbit
observations. A law obtained from the same sources and readings on
those sampled phase tubes is unchanged as well. It is not a claim
that an arbitrary source illuminating additional trajectories would
remain unchanged, and it is not a construction inside every billiard
subclass. The theorem identifies a specific unsampled-orbit freedom;
it does not assert that each of the sufficient coverage hypotheses
in Theorem~\ref{thm:v83-return} is separately necessary.

\subsection{Physical common clocks}

\begin{proposition}[One source and common deadlines]
\label{prop:v83-physical}
Suppose the actual actions on finitely many regular gates satisfy
$0\le W_b\le H$. For $J=q+3$, choose $h>0$, an independent
uniform release delay on $[0,D]$, $D=H+(J+1)h$, and deadlines
$T_j=H+jh$. Use one normalized phase source at every setting.
After the time gate, apply branchwise retention with an unknown
polynomial logarithm of degree at most $q$ on this finite clock
interval, and two conditionally independent readings whose laws do
not depend on the deadline. Then the observations have the form
\eqref{eq:v83-model}, with residual and unused-delay margins at
least $h$. No component audit is needed to choose these deadlines.
\end{proposition}
\begin{proof}
The delay event $\alpha+W_b\le T_j$ has probability
$(T_j-W_b)/D$. Its residual lies in $[h,D-h]$ for every component.
Writing $S_b$ for a regular generating sheet, the change from initial
momentum to terminal position gives the successful density
\[
 \frac{\rho(x,-S_{b,x})|\det S_{b,xy}|}{D}
     \eta_b\exp\left(\sum_{\nu=1}^q\theta_{b,\nu}T_j^\nu\right)
     (T_j-W_b).
\]
All factors outside the exponential and residual are independent of
$j$ and are absorbed in $a_b$. On a finite interval the baseline
$\eta_b$ can be chosen strictly positive and sufficiently small that
the retention lies strictly below one. Independent randomizations
conditional on the trajectory give the product channel columns.
A branch-blind detector or unknown exposure only rescales the matrix.
\end{proof}

For $q=1$ a survival detector with phase-dependent rate is one such
mechanism. The two independent readings may be produced by separate
randomized sensor stages conditional on the trajectory. Their joint
conditional independence and deadline invariance are apparatus
assumptions, not consequences of copying a noisy observation. The
coefficients are unknown within the prescribed finite-dimensional
family; no unrestricted detector is being called identifiable.

\subsection{A nonconvex return system and its persistence}

\begin{proposition}[Three sheets through an action crossing]
\label{prop:v83-shear}
There is a nonconvex exact return system with three regular sheets,
two of which have coincident actions along a curve, satisfying the
log-affine hypotheses of Theorem~\ref{thm:v83-return}. These properties
persist on a sufficiently small smooth neighborhood of the lift,
on slightly smaller fixed phase and endpoint domains.
\end{proposition}
\begin{proof}
On $\R^2$ set
\[
 \phi(p)=p^3-p,\quad G(p)=\tfrac14p^4-\tfrac12p^2,\quad
 F(x,p)=(x+\phi(p),p),\quad A(p)=1+\tfrac34p^4-\tfrac12p^2.
\]
Then $F^*\lambda-\lambda=p\phi'(p)\dd p=dA$ and
$A\ge11/12$. The positive-roof quotient preserves $dt+p\,dx$
and has this lift as its return system. For count $n$, put
$t=(y-x)/n$ and let $p_-(t),p_0(t),p_+(t)$ solve $\phi(p)=t$
near $-1,0,1$. On a small interval about zero they are smooth. The
absolute sheets are
\[
 S_{n,b}(x,y)=n+p_b(t)(y-x)-nG(p_b(t))=nA(p_b(t)).
\]
They satisfy $-S_{n,b,x}=S_{n,b,y}=p_b(t)$ and
$S_{n,b,xy}=-1/(n\phi'(p_b(t)))\ne0$.
On the diagonal the outer actions equal $5n/4$, whereas the central
action is $n$. The derivative of the outer action difference in
$y-x$ there is $p_+(0)-p_-(0)=2$. Thus the coincidence is transverse,
while an unequal-action anchor persists throughout a small compact gate.
The central and outer mixed derivatives have opposite signs.

Choose a detector survival rate $r(p)=r_0+r_1p>0$ on the compact
momentum patches, with $r_1\ne0$, and retention
$\eta\exp[-r(p)(T-T_1)]$, $0<\eta<1$. The rates separate the
outer components at the action crossing. For one reading use
\[
 (1,e^{\beta p},e^{2\beta p})/(1+e^{\beta p}+e^{2\beta p}),
                       \qquad \beta\ne0,
\]
and for the other use the same construction with a different unknown
nonzero coefficient. Draw the readings independently conditional on
phase. Their three columns are normalized Vandermonde columns at
distinct momenta, so they have full rank and strictly positive entries.
The observer is not given the momenta, response coefficients, or rates.
A fixed smooth normalized source positive on the compact patches,
combined with Proposition~\ref{prop:v83-physical}, supplies the clocks.

For counts two and three the momentum is unchanged and the horizontal
shifts are bounded on these patches. Choose finitely many endpoint
rectangles for the input and output neighborhoods in
$g=g^3(g^2)^{-1}$ and then smaller target domains with fixed collars.
This supplies the cover in Theorem~\ref{thm:v83-return}.

Finally take a sufficiently small exact smooth perturbation on a
neighborhood of the finite orbit tubes just chosen. The regular
endpoint equations for both iterates continue their three branches by
the implicit function theorem; no other branch is included in this
restricted compact source class. Their mixed-Jacobian, channel-rank,
residual, and anchor margins persist. The nonzero derivative of the
outer action difference continues a crossing curve in smaller gates.
Keep the sensor functions of initial momentum and their independent
randomizations fixed; their rank and rate separation persist as well.
Collars allow the same finite gate families to cover slightly smaller
phase targets. Positive roofs persist by the strict lower bound.
This proves the neighborhood assertion, rather than only the single
explicit example. Fold points $\phi'(p)=0$ are outside these domains.
\end{proof}

\section{Exact boundaries of the observation theorem}
\label{sec:v83-boundaries}

\begin{proposition}[Equal actions and unrestricted detectors]
\label{prop:v83-obstructions}
If all components have the same action $w$, projective observations
at arbitrarily many deadlines do not determine $w$. If the branchwise
detector may vary freely with deadline, small changes of unequal
actions can also be hidden by an exactly compensating detector.
\end{proposition}
\begin{proof}
In the first case the common factor $T_j-w$ cancels from every
normalized matrix. In the second case, replace an admissible action
$W_b$ by a nearby $W'_b<T_1$ and replace its detector on the fixed
clock interval by
\[
             e'_b(T)=e_b(T)\frac{T-W_b}{T-W'_b}.
\]
Every unnormalized component weight is unchanged. Strict residual and
retention margins keep this a valid probability for sufficiently small
changes. The compensating logarithm need not belong to the prescribed
polynomial family. Thus this construction does not contradict the sharp
positive theorem.
\end{proof}

\begin{proposition}[One uncalibrated reading does not suffice]
\label{prop:v83-one-view}
Even with two unequal actions, strict probability margins, and a
full-rank single readout, projective observations from one unknown
reading need not determine the actions at any number of deadlines.
\end{proposition}
\begin{proof}
Take unit weights and zero detector coefficients. In the first model
use actions $(0,1)$ and channel columns $(.2,.8)^{\mathsf T}$,
$(.8,.2)^{\mathsf T}$. In the second use actions $(-1,2)$ and
columns $(.4,.6)^{\mathsf T}$, $(.6,.4)^{\mathsf T}$.
For every $T>2$ the unnormalized single-reading vector is, in both
models,
\[
                       (T-.8,T-.2)^{\mathsf T}.
\]
The two channel matrices are positive and invertible, but the unordered
action sets differ. Their normalized laws, and even these raw vectors,
are identical. Common scaling makes any finite-clock success masses
valid probabilities without altering this equality.
\end{proof}

The obstruction is to one unknown channel over the full stated class;
it does not rule out identification with a calibrated sensor or other
additional information. Together with the clock lower bound, it locates
three distinct requirements: enough clock contrasts, enough independent
readout information, and an unequal-action anchor. Phase coverage is a
separate geometric requirement. None of these distinctions changes the
absolute-action conclusion in the identifiable class.
